\documentclass[11pt, a4paper]{article}

\usepackage[a4paper, top=2.8cm, bottom=2.8cm, left=2.5cm, right=2.5cm, headheight=16pt]{geometry}
\usepackage{fontspec}
\usepackage[german, bidi=basic, provide=*]{babel}

\babelprovide[import, onchar=ids fonts]{german}
\babelprovide[import, onchar=ids fonts]{english}

\babelfont{rm}{Noto Sans}
\babelfont[german]{rm}{Noto Sans}

\usepackage{xcolor}
\usepackage{titlesec}
\usepackage{tcolorbox}
\tcbuselibrary{listings}
\usepackage{tabularx}
\usepackage{fancyhdr}
\usepackage[colorlinks=true, linkcolor=eopurple, urlcolor=eocyan]{hyperref}
\usepackage{enumitem}
\usepackage{booktabs}

% EO Corporate Design Farben
\definecolor{eopurple}{RGB}{139, 92, 246}
\definecolor{eocyan}{RGB}{6, 182, 212}
\definecolor{eodark}{RGB}{15, 23, 42}
\definecolor{codebg}{RGB}{241, 245, 249}
\definecolor{codeframe}{RGB}{203, 213, 225}

% Titel-Formatierung
\titleformat{\section}{\color{eopurple}\normalfont\Large\bfseries}{\thesection}{1em}{}
\titleformat{\subsection}{\color{eocyan}\normalfont\large\bfseries}{\thesubsection}{1em}{}
\titleformat{\subsubsection}{\color{eodark}\normalfont\normalsize\bfseries}{\thesubsubsection}{1em}{}

% Code-Block Styling
\lstdefinestyle{eo_code}{
    basicstyle=\small\ttfamily,
    breaklines=true,
    breakatwhitespace=true,
    tabsize=2,
    showstringspaces=false
}

\newtcblisting{codeblock}{
    listing only,
    listing style=eo_code,
    colback=codebg,
    colframe=codeframe,
    arc=4pt,
    boxrule=0.5pt,
    left=4pt,
    right=4pt,
    top=4pt,
    bottom=4pt
}

\newcommand{\code}[1]{\colorbox{codebg}{\texttt{\small\detokenize{#1}}}}

% Änderungs-Marker v1.0
\newcommand{\changed}[1]{\textcolor{eopurple}{\textbf{[v1.0]}}~#1}
\newcommand{\changedtitle}[1]{\textcolor{eopurple}{#1}}

% Kopf- und Fußzeile
\pagestyle{fancy}
\fancyhf{}
\fancyhead[L]{\textbf{\textcolor{eopurple}{EO} \textcolor{eocyan}{einfach-online.dev}}}
\fancyhead[R]{\small Carp24 \& Supabase — Administrator Playbook (v1.0)}
\fancyfoot[C]{\thepage}
\fancyfoot[R]{\scriptsize EO-ADMIN-2026-CARP24-01-REV}

\hyphenation{Su-pa-base Car-p24 Post-gres Go-Tru-e Post-gREST Supa-base Ad-mi-nis-tra-ti-on}

\begin{document}

\begin{center}
    \vspace*{1.0cm}
    {\color{eopurple}\Huge\bfseries Carp24 \& Supabase}\\[0.4cm]
    {\color{eocyan}\Large\bfseries Administrator Playbook}\\[0.3cm]
    {\large Einfach verständlich: Was ist Supabase, was kann es, und wie verwaltest du Carp24 als Admin.}\\[0.8cm]
    {\normalsize Version 1.0 \quad\textbar\quad 21.08.2026 \quad\textbar\quad Dokument-ID: EO-ADMIN-2026-CARP24-01-REV}\\[0.3cm]
    {\normalsize einfach-online.dev \quad\textbar\quad Philipp Schlemmer}
    \vspace*{1.0cm}
\end{center}

\tableofcontents
\newpage

% ============================== CHANGELOG ==============================
\section{Changelog \& Änderungsverzeichnis}

\begin{tabularx}{\textwidth}{llX}
\toprule
\textbf{Version} & \textbf{Datum} & \textbf{Änderung} \\
\midrule
v1.0 & 21.08.2026 & Initiales Dokument: Supabase-Grundlagen, Carp24-Features, Admin-Dashboard v2, Einstellungen, Betrieb \& Fehlerbehandlung \\
\bottomrule
\end{tabularx}

\textbf{Hinweis:} Dieses Dokument enthält \textbf{keine Passwörter oder Secret-Werte}. Alle Schlüssel liegen im Passwort-/Secret-Handling des Orchestrators (Hermes) und in den lokalen \code{.env}-Dateien des Servers.

\newpage

% ============================== TEIL A: SUPABASE ==============================
\section{Teil A — Supabase: Was ist das überhaupt?}

\subsection{Was ist Supabase?}
Supabase ist eine \textbf{Open-Source-Alternative zu Firebase} (Google). Es ist eine Komplett-Plattform für das Backend einer Web-App und besteht im Kern aus:

\begin{itemize}
    \item \textbf{PostgreSQL-Datenbank} — eine der mächtigsten relationalen Datenbanken (SQL)
    \item \textbf{Authentifizierung (Auth)} — Registrierung, Login, Passwort-Reset, Social-Login (Google/Facebook), Zwei-Faktor (MFA)
    \item \textbf{Auto-API} — Jede Tabelle bekommt automatisch eine REST-API (PostgREST); kein Backend-Code nötig
    \item \textbf{Storage} — Datei-Upload (z.B. Fang-Fotos), S3-kompatibel
    \item \textbf{Realtime} — Live-Updates im Browser (WebSocket), z.B. für den Chat
    \item \textbf{Studio} — eine Weboberfläche, in der du Tabellen, Daten und Einstellungen direkt ansehen kannst
\end{itemize}

Supabase läuft bei uns \textbf{selbst gehostet} auf dem Server als Docker-Container — nicht in der Cloud von Supabase. Das heißt: Alle Daten bleiben auf unserer Infrastruktur (DSGVO-freundlich).

\subsection{Die Bausteine (Container) auf unserem Server}

\begin{tabularx}{\textwidth}{llX}
\toprule
\textbf{Container} & \textbf{Aufgabe} & \textbf{Hinweis} \\
\midrule
supabase-db & PostgreSQL 17 & Die eigentliche Datenbank (Port 5432 intern) \\
supabase-auth & GoTrue v2.189 & Login/Registrierung/Passwort/MFA/JWT \\
supabase-rest & PostgREST v14 & Auto-REST-API (Port 3000 intern) \\
supabase-storage & Storage v1.60.4 & Datei-Upload (Fang-Fotos) \\
supabase-realtime & Realtime v2.102 & Live-WebSockets (Chat) \\
supabase-kong & Kong 3.9.3 & API-Gateway: zentrale Tür (extern 100.93.250.103:8055) \\
supabase-studio & Studio & Web-UI zur Datenbank-Verwaltung \\
supabase-meta & postgres-meta & Verwaltungs-API hinter Studio \\
supabase-pooler & Supavisor & Verbindungs-Pooling (Port 6543) \\
\bottomrule
\end{tabularx}

\textbf{Wichtig für dich:} Alle Wege nach außen laufen über \textbf{Kong} (Port 8055). Die Carp24-Web-App erreicht Supabase über den Apache-Proxy \code{:8094}, der auf Kong weiterleitet (\code{/auth}, \code{/rest}, \code{/storage}, \code{/realtime}).

\subsection{Kernkonzepte, die du als Admin verstehen solltest}

\begin{description}
    \item[Tabellen] Daten liegen in Tabellen (z.B. \code{catches} für Fänge). Jede Zeile = ein Datensatz.
    \item[RLS (Row Level Security)] Eine Sicherheits-Ebene \textbf{in der Datenbank}: Legt fest, wer welche Zeilen sehen/ändern darf (z.B. „jeder nur seine eigenen Fänge"). Ohne RLS wäre alles für jeden erreichbar, der den API-Schlüssel kennt.
    \item[Policies] Die konkreten RLS-Regeln pro Tabelle (z.B. \code{USING (user\_id = auth.uid())}).
    \item[Migrationen] SQL-Dateien, die Schema-Änderungen dokumentieren (Ordner \code{supabase/migrations/}). Sie sind der „Stammbaum" der Datenbank — niemals einfach so ändern, sondern neue Datei anlegen.
    \item[Buckets] Speicher-Container für Dateien (bei Carp24: \code{catch-photos}).
    \item[Publicationen] Listen von Tabellen, die Realtime-Updates senden dürfen (z.B. \code{chat\_messages}).
    \item[RPC-Funktionen] Eigene SQL-Funktionen, die von der App aufgerufen werden (z.B. \code{import\_catches} für den CSV-Import).
    \item[JWT / Session] Nach dem Login bekommt der Browser ein signiertes Token (JWT). Die API erkennt daran den Benutzer (\code{auth.uid()}).
\end{description}

\subsection{Wie kommst du an die Datenbank?}

Es gibt drei Wege:
\begin{enumerate}
    \item \textbf{Supabase Studio} (komfortabel, Web-UI) — Tabellen, Zeilen, SQL-Editor
    \item \textbf{psql} (Kommandozeile, direkt im Container) — für Migrationen und SQL
    \item \textbf{REST-API} über Kong — was die App selbst nutzt
\end{enumerate}

Migrationen einspielen (Standardweg, vom Orchestrator ausgeführt):
\begin{codeblock}
cd /mnt/projekte/carp24-fangbuch
PW=$(docker exec supabase-db printenv PGPASSWORD)
docker exec -i --env PGPASSWORD="$PW" supabase-db psql -h localhost -U supabase_admin -d postgres < supabase/migrations/00XX_name.sql
\end{codeblock}

\newpage
% ============================== TEIL B: CARP24 ==============================
\section{Teil B — Carp24: Die Plattform}

\subsection{Was ist Carp24?}
Carp24 ist eine \textbf{Karpfenangler-Community mit digitalem Fangbuch} (eigenes Produkt von einfach-online.dev). Angler erfassen ihre Fänge (Gewicht, Fischart, Gewässer, Foto, Wetter), sammeln Badges, sehen Jahres-Rückblicke und tauschen sich in Forum, Chat, Community-Board und Marktplatz aus. Es gibt ein Free-/Pro-Modell (Stripe) und einen KI-Angelassistenten.

\textbf{Technik:}
\begin{itemize}
    \item Frontend: \textbf{Astro 7} (SSR, Node-Adapter), Code in \code{/mnt/projekte/carp24-fangbuch/web/}
    \item Backend: \textbf{Supabase self-hosted} (siehe Teil A)
    \item Dev-Lauf: \textbf{http://100.93.250.103:8094} (Apache → Node :4321)
    \item Deploy: systemd-Dienst \code{carp24-web} (\code{sudo systemctl restart carp24-web})
    \item Branch: \code{sprint-1}, Migrationen 0013–0030
\end{itemize}

\subsection{Alle Produkt-Features (Stand v1.0)}

\begin{tabularx}{\textwidth}{llX}
\toprule
\textbf{Bereich} & \textbf{Feature} & \textbf{Kurzbeschreibung} \\
\midrule
Kern & Fänge erfassen & Fang-Formular mit Foto, Gewicht, Art, Gewässer, Wetter, Koordinaten \\
Kern & Statistik & Diagramme \& Auswertungen der eigenen Fänge \\
Kern & Offline-Modus & Fänge ohne Internet zwischenspeichern \\
A & CSV-Import & Alte Fanglisten als CSV importieren (max. 500 Zeilen) \\
A & Badges & 10 Auszeichnungen (Erster Fang, 100 Kilo, Regenfischer \dots) \\
A & Jahres-Rückblick & KPI-Karten, Monats-Balken, Mond-Statistik, Teilen \\
A & Trips & Angelausflüge anlegen und Fänge zuordnen \\
B & OAuth-Login & Google-/Facebook-Buttons (Creds konfigurierbar) \\
C & Community-Board & Öffentliche Top-Fänge teilen (Public-Toggle) \\
C & Forum & Threads + Antworten mit 4 Kategorien, Melden, Moderation \\
C & Chat & Realtime-Chat mit Live-WebSocket, Melden, Löschen \\
D & Pro-Gating & Free vs. Pro (Locks), Admin sieht alles \\
D & Stripe-Abo & Checkout + Webhook setzt \code{is\_pro} (Test-Modus) \\
E & i18n & Deutsch/Englisch-Umschalter, Rechtstexte zweisprachig \\
E & Web-Push & Benachrichtigungen (VAPID), Antwort-Push, Tages-Erinnerung \\
E & KI-Assistent & Chat mit OpenRouter-API, keine Speicherung (DSGVO) \\
E & Marktplatz & Anzeigen (Foto, Preis), In-App-Kontakt-Nachrichten \\
X & DSGVO-Export & Eigene Daten exportieren (Art. 15) \\
X & DSGVO-Löschung & Eigene Daten löschen (Art. 17) \\
X & Meldungen & \code{content\_reports} — User melden Inhalte (DSA) \\
\bottomrule
\end{tabularx}

\subsection{Das Admin-Dashboard (dein wichtigstes Werkzeug)}

Erreichbar über \code{http://100.93.250.103:8094/admin/} nach Login mit einem Konto mit Rolle \code{ADMIN} oder \code{MODERATOR} (Header-Link \textbf{ADMIN}). Sechs Bereiche:

\subsubsection{Übersicht}
Zahlen auf einen Blick: Nutzer, Fänge, offene Meldungen, Threads, Chat-Nachrichten, Anzeigen, Trips, aktive Pro-User, Gebannte. Dazu ein 14-Tage-Wachstums-Diagramm (Nutzer/Fänge) und eine Integrations-Status-Karte (welche Dienste konfiguriert sind: Stripe, KI, Push, Service-Role — nur Ja/Nein, keine Schlüssel).

\subsubsection{Reports (Meldungen moderieren)}
Nutzer können Inhalte melden (Fang, Forum-Thread, Forum-Antwort, Chat-Nachricht, Marktplatz-Anzeige). Hier siehst du:
\begin{itemize}
    \item Offene/erledigte Meldungen (Tabs), mit Grund, Melder und Ziel-Inhalt im Kontext
    \item Button \textbf{ERLEDIGT} — Meldung abschließen
    \item Button \textbf{INHALT LÖSCHEN} — Ziel-Inhalt entfernen und Meldung erledigen
    \item Button \textbf{BANNEN} — wenn ein User mehrfach auffällig ist (≥3 offene Meldungen)
\end{itemize}

\subsubsection{Nutzer verwalten}
\begin{itemize}
    \item Suche (Name/E-Mail) + Filter (Rolle, Status, Pro)
    \item \textbf{Rolle} ändern: USER / MODERATOR / ADMIN
    \item \textbf{Pro} an/aus (ohne Stripe-Zahlung manuell setzen)
    \item \textbf{BANNEN / ENTBANNEN} (mit Grund) — blockiert auch den Login (GoTrue-Ban)
    \item \textbf{PASSWORT ZURÜCKSETZEN} — erzeugt ein temporäres Passwort (Support-Fall)
    \item \textbf{DETAIL} — öffnet den User-Feed: alle Fänge, Forum-Beiträge, Chat-Nachrichten, Anzeigen, Trips, Badges und Meldungen des Users auf einen Blick; Inhalte direkt löschbar
\end{itemize}

\subsubsection{Inhalte moderieren}
Tabs für Fänge, Forum, Chat und Marktplatz. Durchsuchen und einzelne Inhalte löschen (Fänge werden „weich" gelöscht / soft-delete, Rest hart).

\subsubsection{Revenue (Abo-Übersicht)}
Aktive Pro-User, Free-User, Pro-Quote, Liste der Pro-User, Stripe-Events-Status. Hinweis, wenn Stripe noch nicht konfiguriert ist.

\subsubsection{Audit (Protokoll)}
Wer hat wann was getan: Rollen-Änderungen, Banns, Passwort-Resets, Report-Bearbeitung, Inhalts-Löschung. Wichtig für die DSGVO-Nachweispflicht (Accountability).

\subsection{Die wichtigsten Datenbank-Tabellen}

\begin{tabularx}{\textwidth}{llX}
\toprule
\textbf{Tabelle} & \textbf{Migration} & \textbf{Inhalt} \\
\midrule
profiles & Basis & Nutzerprofile (display\_name, role, is\_pro, banned\_at, \dots) \\
catches & Basis & Fänge (species, weight\_kg, water\_name, photo, is\_public) \\
trips & 0019 & Angelausflüge (catches.trip\_id verknüpft) \\
badges / user\_badges & 0018 & 10 Badges + Verleihungen \\
forum\_threads / forum\_posts & 0021 & Forum (category ab 0028) \\
chat\_messages & 0022 & Realtime-Chat \\
content\_reports & 0017 & Meldungen (status open/resolved) \\
marketplace\_items & 0026 & Marktplatz-Anzeigen \\
marketplace\_contacts & 0029 & In-App-Kontakt-Nachrichten \\
push\_subscriptions & 0025 & Push-Abos der Browser \\
stripe\_events & 0024 & Stripe-Webhook-Ereignisse (Idempotenz) \\
admin\_audit\_log & 0030 & Admin-Aktionen (Audit) \\
\bottomrule
\end{tabularx}

\newpage
% ============================== TEIL C: EINSTELLUNGEN ==============================
\section{Teil C — Wo kann man was einstellen?}

\subsection{Konfigurationsdateien auf dem Server}

\textbf{1. \code{/mnt/projekte/carp24-fangbuch/web/.env} — Einstellungen der Web-App:}

\begin{tabularx}{\textwidth}{lXX}
\toprule
\textbf{Variable} & \textbf{Wofür?} & \textbf{Status (21.08.)} \\
\midrule
PUBLIC\_SUPABASE\_URL & Basis-URL zu Supabase (über Apache :8094) & gesetzt \\
PUBLIC\_SUPABASE\_ANON\_KEY & Öffentlicher API-Schlüssel (Browser) & gesetzt \\
SUPABASE\_SERVICE\_ROLE\_KEY & Admin-Schlüssel (nur Server!) & gesetzt \\
STRIPE\_SECRET\_KEY & Stripe-Test-Server-Schlüssel & \textbf{wartet auf Philipp} \\
STRIPE\_WEBHOOK\_SECRET & Stripe-Webhook-Signatur & \textbf{wartet auf Philipp} \\
STRIPE\_PRICE\_ID & Stripe-Preis-ID für Pro-Abo & \textbf{wartet auf Philipp} \\
OPENROUTER\_API\_KEY & KI-Assistent (OpenRouter) & \textbf{wartet auf Philipp} \\
VAPID\_PUBLIC\_KEY / VAPID\_PRIVATE\_KEY / VAPID\_SUBJECT & Web-Push & gesetzt \\
CRON\_KEY & Schutz für den Erinnerungs-Cron & \textbf{noch zu setzen} \\
\bottomrule
\end{tabularx}

\textbf{2. \code{/mnt/projekte/carp24-fangbuch/infra/.env} — Supabase-Stack (nur Orchestrator ändern):}
Enthält JWT-Secret, ANON\_KEY, SERVICE\_ROLE\_KEY, API-URLs und die GoTrue-Provider-Platzhalter für Google/Facebook-Login (\code{GOTRUE\_EXTERNAL\_GOOGLE\_ENABLED} etc.).

\textbf{3. Supabase Studio} — für Tabellen, Zeilen, SQL und RLS direkt. Nicht öffentlich erreichbar; Zugriff über Docker/SSH.

\subsection{Google-/Facebook-Login aktivieren (OAuth)}

\begin{enumerate}
    \item In Google Cloud Console / Facebook Developers eine OAuth-App anlegen
    \item Redirect-URL: \code{http://100.93.250.103:8094/auth/v1/callback}
    \item Client-ID + Secret in \code{infra/.env} eintragen (\code{GOTRUE\_EXTERNAL\_GOOGLE\_CLIENT\_ID}, \code{...\_SECRET}, FB analog)
    \item Auth-Container neu starten: \code{docker compose restart supabase-auth}
    \item Test: Login-Seite → „Mit Google anmelden"
\end{enumerate}

\subsection{Stripe aktivieren (Pro-Abo)}

\begin{enumerate}
    \item Stripe-Konto (Test-Modus) → API-Keys: \code{STRIPE\_SECRET\_KEY} (sk\_test\dots) in \code{web/.env}
    \item Produkt + Preis anlegen → \code{STRIPE\_PRICE\_ID}
    \item Webhook einrichten: URL \code{https://<domain>/api/stripe/webhook}, Event \code{checkout.session.completed} + \code{customer.subscription.deleted/updated} → \code{STRIPE\_WEBHOOK\_SECRET}
    \item \code{SUPABASE\_SERVICE\_ROLE\_KEY} muss gesetzt sein (ist es)
    \item Test: Premium-Seite → „JETZT UPGRADEN" → Test-Checkout → \code{is\_pro=true}
\end{enumerate}

\subsection{KI-Assistent aktivieren}
OpenRouter-Key (\code{OPENROUTER\_API\_KEY}) in \code{web/.env} setzen. Ohne Key liefert der Assistent sauber einen 503-Hinweis (Feature ist dann deaktiviert statt kaputt).

\subsection{Web-Push-Erinnerung (täglicher Cron)}
Die API-Route \code{/api/push/daily-reminder} wird per systemd-Timer (täglich 08:00) aufgerufen und braucht den Header \code{x-cron-key}. Vorher \code{CRON\_KEY} in \code{web/.env} setzen, dann richtet der Orchestrator den Timer ein.

\subsection{Pro-Status \& Rollen manuell setzen}
Im Admin-Dashboard unter \textbf{Nutzer}: Rolle-Select (USER/MODERATOR/ADMIN) und Pro-Checkbox. Alternativ direkt in der Datenbank:
\begin{codeblock}
docker exec --env PGPASSWORD="$PW" supabase-db psql -h localhost -U supabase_admin -d postgres \
  -c "UPDATE profiles SET role='ADMIN', is_pro=true WHERE id='<user-id>';"
\end{codeblock}

\newpage
% ============================== TEIL D: BETRIEB ==============================
\section{Teil D — Betrieb, Wartung \& Fehler}

\subsection{Neue Version deployen}

\begin{enumerate}
    \item Code ändern lassen (OpenCode-Briefings, nie direkt am Server)
    \item Build: \code{cd /mnt/projekte/carp24-fangbuch/web \&\& npm run build}
    \item Neustart: \code{sudo systemctl restart carp24-web}
    \item Prüfen: \code{systemctl is-active carp24-web} → \code{active}; Seite per Browser testen
\end{enumerate}

\textbf{Hinweis:} Während eines Builds kann der laufende Server kurz 500er liefern (dist wird überschrieben). Nach dem Restart ist alles wieder normal. Nicht während eines Builds testen.

\subsection{Migrationen einspielen}
Neue SQL-Datei liegt in \code{supabase/migrations/}. Einspielen siehe Teil A (psql-Kommando). Nach FK-/Schema-Änderungen ggf. \code{docker compose restart supabase-rest} (Schema-Reload).

\subsection{Backups \& Wiederherstellung}
Backup-Lauf täglich 02:00 (Docker-Dumps). Restore-Drill dokumentiert im Vault (RESTORE.md). Vor größeren Änderungen: Backup prüfen.

\subsection{Logs lesen}
\begin{codeblock}
sudo journalctl -u carp24-web --no-pager -n 100        # Web-App-Log
sudo journalctl -u carp24-web --since "10 min ago"     # aktuelle Fehler
docker logs supabase-auth --tail 50                    # Auth-Log
\end{codeblock}

\subsection{Typische Fehler und ihre Lösungen}

\begin{tabularx}{\textwidth}{lX}
\toprule
\textbf{Symptom} & \textbf{Ursache / Fix} \\
\midrule
Fangliste zeigt „keine Fänge" & Null-Filter falsch: IMMER \code{.is("deleted\_at", null)} verwenden, nie \code{.eq(..., null)} \\
500 „supabaseKey is required" & SERVICE\_ROLE\_KEY in web/.env fehlt oder Build lief vor dem Setzen → Key setzen + \textbf{neu bauen} \\
Forum-Seite 500 „t2 is not a function" & i18n-Überschattung: In \code{map}-Callbacks nie den Parameter \code{t} nennen \\
Foto-Upload schlägt fehl & Storage-Bucket \code{catch-photos} + \code{GLOBAL\_S3\_BUCKET} muss gesetzt sein \\
Login/Realtime funktioniert nicht & Kong muss auf erreichbarer IP binden; Apache-Proxy \code{/auth} \code{/realtime} mit WebSocket-Upgrade \\
Seite crasht client-seitig & PageErrors im Browser prüfen; anonyme Besucher müssen auf \code{/login} redirectet werden (SSR-Guard) \\
\bottomrule
\end{tabularx}

\subsection{Sicherheitsregeln (unbedingt)}

\begin{itemize}
    \item \textbf{Niemals Secrets ins Git-Repo} — \code{.env}-Dateien sind ausgeschlossen (gitignore). Secret-Scan blockt automatisch
    \item \code{SUPABASE\_SERVICE\_ROLE\_KEY} nur auf dem Server (nie im Browser-Code)
    \item RLS ist die Sicherheitsgrenze — Policies nie ohne Review ändern
    \item Admin-Aktionen landen im \code{admin\_audit\_log} (DSGVO)
    \item Benachrichtigungs-/E-Mail-Tests nur mit echten Adressen (keine Fake-Domains)
\end{itemize}

\subsection{Wichtige technische Merksätze}
\begin{itemize}
    \item DB heißt \code{postgres}; User \code{supabase\_auth\_admin} für Auth-Daten, \code{supabase\_admin} für DDL
    \item Supabase-Stack NICHT einzeln pinnen — gesamten Stack auf einem Release-Stand halten (Doku-Regel MASTER\_SPRINT)
    \item Realtime-Chat braucht die Tabelle in der Publication \code{supabase\_realtime}
    \item Temp-Watchdog (systemd) räumt hängende Test-Browser automatisch auf
\end{itemize}

\newpage
% ============================== GLOSSAR ==============================
\section{Glossar}

\begin{description}
    \item[Admin-Dashboard] Die interne Verwaltungsoberfläche von Carp24 (\code{/admin/})
    \item[Badge] Spielerische Auszeichnung im Profil (10 Stück)
    \item[DSGVO] Datenschutz-Grundverordnung (EU); Art. 15/17 = Auskunft/Löschung
    \item[DSA] Digital Services Act (EU); Meldepflichten für User-Content
    \item[GoTrue] Supabase-Auth-Dienst (Login, JWT, MFA)
    \item[is\_pro] Flag auf dem Profil: zahlender Pro-User
    \item[JWT] Signiertes Login-Token (JSON Web Token)
    \item[Migration] Versionierte SQL-Änderung an der Datenbank
    \item[PostgREST] Werkzeug, das aus SQL-Tabellen automatisch eine REST-API macht
    \item[RLS] Row Level Security: Zeilen-Sicherheit in der Datenbank
    \item[RPC] Eigene SQL-Funktion, die die App aufruft
    \item[VAPID] Schlüsselpaar für Web-Push-Benachrichtigungen
\end{description}

\end{document}
